% !TeX program = xelatex
\documentclass[11pt,a4paper]{article}

% Компилировать с XeLaTeX или LuaLaTeX.
\usepackage[a4paper,margin=0.62in]{geometry}
\usepackage{fontspec}
\IfFontExistsTF{Lato}{\setmainfont{Lato}}{\setmainfont{Noto Sans}}
\usepackage{polyglossia}
\setdefaultlanguage{russian}
\usepackage{microtype}
\usepackage{tabularx}
\usepackage{array}
\usepackage{enumitem}
\usepackage{xcolor}
\usepackage{hyperref}
\usepackage{fancyhdr}
\usepackage{lastpage}

\definecolor{linkblue}{RGB}{20,33,110}
\hypersetup{
  colorlinks=true,
  urlcolor=linkblue,
  linkcolor=linkblue,
  citecolor=linkblue
}

\pagestyle{fancy}
\fancyhf{}
\renewcommand{\headrulewidth}{0pt}
\renewcommand{\footrulewidth}{0pt}
\cfoot{\small \thepage\ из \pageref{LastPage}}

\setlength{\parindent}{0pt}
\setlength{\parskip}{0pt}
\setlength{\tabcolsep}{0pt}
\setlist[itemize]{leftmargin=1.15em,itemsep=0.08em,topsep=0.15em,parsep=0pt,partopsep=0pt}

\newcolumntype{L}{>{\raggedright\arraybackslash}p{0.145\textwidth}}
\newcolumntype{R}{>{\raggedright\arraybackslash}p{0.835\textwidth}}

\newcommand{\cvrow}[2]{%
  \noindent
  \begin{minipage}[t]{0.145\textwidth}
    \vspace{0pt}%
    \raggedright
    \fontsize{8.5}{10}\selectfont
    \bfseries\MakeUppercase{#1}
  \end{minipage}%
  \hfill
  \begin{minipage}[t]{0.835\textwidth}
    \vspace{0pt}%
    #2
  \end{minipage}\par
}

\newcommand{\entry}[2]{%
  \begin{tabularx}{\linewidth}{@{}X >{\raggedleft\arraybackslash}p{0.28\linewidth}@{}}
    #1 & #2
  \end{tabularx}
}

\newcommand{\pub}[2]{%
  #1 #2\par
  \vspace{0.28em}
}

\newcommand{\eduentry}[2]{%
  \noindent
  \begin{tabularx}{\linewidth}
    {@{}X@{\hspace{0.8em}}
    >{\raggedleft\arraybackslash}p{0.30\linewidth}@{}}
    #1 & \small #2
  \end{tabularx}\par
  \vspace{0.15em}
}

\newcommand{\edudetail}[1]{%
  \noindent\hspace*{1.15em}$\diamond$ #1\par
}

\begin{document}

{\fontsize{24}{28}\selectfont\bfseries Артур Ясеновец}\par
\vspace{0.18em}
\hrule
\vspace{0.65em}

\cvrow{Текущая\\позиция}{%
  \begin{tabularx}{\linewidth}[t]
    {@{}>{\raggedright\arraybackslash}X
    @{\hspace{0.45em}}
    >{\raggedright\arraybackslash}p{0.46\linewidth}@{}}

    \begin{tabular}[t]{@{}l@{}}
      \textbf{Магистрант-исследователь}\\[-0.15em]
      \href{https://www.cqupt.edu.cn/}
        {Чунцинский ун-т почты и телеком.}
    \end{tabular}
    &
    \begin{tabular}[t]{@{}l@{}}
      \textit{Эл. почта:}
      \href{mailto:L202420002@stu.cqupt.edu.cn}{L202420002@stu.cqupt.edu.cn}\\[-0.15em]
      \textit{Список публикаций:}
      \href{https://orcid.org/0009-0008-0008-3746}{ORCID}
    \end{tabular}

  \end{tabularx}
}

\vspace{-0.35em}

\cvrow{Научные\\интересы}{%
  Мои научные интересы лежат в области \textbf{прикладной криптографии} и
  \textbf{технологий повышения конфиденциальности}, с особым вниманием к
  \textbf{гомоморфному шифрованию}, \textbf{приватному извлечению информации},
  \textbf{шифрованию с возможностью поиска} и системам с сохранением конфиденциальности.
  В особенности меня интересуют практические криптографические протоколы для
  конфиденциального поиска по данным и приватного извлечения записей.
}

\vspace{0.75em}

\cvrow{Дополн-ая\\информация}{%
  \textbf{Научные профили:}
  \href{https://iasenovets.github.io/}{Портфолио}
  \;\textbullet\;
  \href{https://linkedin.com/in/iasenovets}{LinkedIn}\par
  \textbf{Рекомендации:} контакты рекомендателей предоставляются по запросу.
}

\vspace{0.75em}

\cvrow{Образование}{%
  \raggedright

  % ── Магистратура ────────────────────────────────
  \textbf{Чунцинский университет почты и телекоммуникаций},
  Чунцин, Китай\par
  \textit{Школа кибербезопасности и информационного права}\par
  \vspace{0.12em}

  \eduentry
    {$\diamond$ \textbf{Магистратура} по направлению «Сетевая и информационная безопасность»}
    {сент. 2024--июль 2027 \textit{(ожидается)}}

  \edudetail{Научный руководитель:
    \href{https://faculty.cqupt.edu.cn/tangfei/en/index.htm}{проф. Тан Фэй}}
  \edudetail{Средний балл: 3,98/4,00}

  \vspace{0.48em}

  % ── Бакалавриат ─────────────────────────────────
  \textbf{Ухтинский государственный технический университет},
  Ухта, Россия\par
  \vspace{0.12em}

  \eduentry
    {$\diamond$ \textbf{Бакалавриат} по направлению «Информационные системы и технологии»}
    {март 2022--июнь 2024}

  \edudetail{Средневзвешенный балл: 4,49/5,00}

  \vspace{0.48em}

  % ── Предыдущее обучение в бакалавриате ─────────
  \textbf{Санкт-Петербургский государственный университет телекоммуникаций
  им. проф. М.\,А. Бонч-Бруевича},
  Санкт-Петербург, Россия\par
  \vspace{0.12em}

  \eduentry
    {$\diamond$ \textbf{Обучение в бакалавриате} по направлению
    «Информационные системы и технологии»}
    {сент. 2020--март 2022}
}

\vspace{0.75em}

\cvrow{Публикации}{%
  \textbf{Опубликованные работы и препринты}\par\vspace{0.25em}

  \pub{\underbar{Artur Iasenovets}, Fei Tang, H. Zhu, P. Wang, and L. Liu.}
  {``\href{https://arxiv.org/abs/2511.02656}
  {Bringing Private Reads to Hyperledger Fabric via Private Information Retrieval},''
  arXiv:2511.02656, 2025. Рукопись находится на рецензировании.}

  \pub{H. Zhu, F. Tang, W. Qin, J. Shan, P. Wang, and \underbar{Artur Iasenovets}.}
  {``\href{https://doi.org/10.1007/s44443-025-00405-8}
  {EVMK-SSE: Efficient and Verifiable Multi-Keyword Symmetric Searchable Encryption},''
  \textit{Journal of King Saud University -- Computer and Information Sciences}, 2025.}

  \pub{\underbar{Artur Iasenovets} and Fei Tang.}
  {``\href{https://doi.org/10.60797/IRJ.2026.163.57}
  {MFCTIIF: Multi-Feed Cyber Threat Intelligence Integration Framework},''
  \textit{International Research Journal}, № 1 (163), 2026.}

  \vspace{0.15em}
  \textbf{Рукопись в подготовке}\par\vspace{0.25em}

  \pub{Piyumi M. S. Rajapaksha$^{*}$, \underbar{Artur Iasenovets$^{*}$},
  Yinxue Yi, and Fei Tang.}
  {``TWAPPA-FRL: Trust-Weighted Privacy-Preserving Aggregation for Federated
  Reinforcement Learning in Collaborative Intrusion Detection,''
  рукопись в подготовке, 2026. $^{*}$Равный вклад.}
}

\vspace{0.75em}

\cvrow{Награды и\\достижения}{%
  \entry
    {\textbf{Стипендия правительства КНР (CSC) для обучения в магистратуре}}
    {2024--2027}
  Полная стипендия на весь период обучения в магистратуре CQUPT.

  \vspace{0.45em}
  \entry
    {\textbf{Победитель конкурса EdTech-проектов с проектом YouKnow}}
    {2023}
  Победа в конкурсе, организованном Физтех-школой МФТИ,
  Фондом развития Физтех-школ и Startech.Base.

  \vspace{0.45em}
  \entry
    {\textbf{Диплом I степени, «Севергеоэкотех»}}
    {2022}
  Первое место в секции «Современные информационные технологии»
  XXIII Международной молодёжной научной конференции.

  \vspace{0.45em}
  \entry
    {\textbf{Региональный демо-день Sber Student Accelerator}}
    {2023}
  Завершил программу Sber Student Accelerator в составе команды AcademAI
  и принял участие в региональном демо-дне Северо-Западного округа.

  \vspace{0.3em}
  \href{https://iasenovets.github.io/files/certificates_selected_en_public_ready_v2.pdf}
       {\small Подтверждающие сертификаты}
}

\vspace*{0.15em}
\cvrow{Научный\\опыт}{%
  \entry
    {\textbf{Магистрант-исследователь} @ Исследовательская группа
    по криптографии и её приложениям, CQUPT}
    {сент. 2024--наст. время}
  \vspace{0.15em}

  \vspace{0.35em}

  \textbf{Приватное извлечение информации для Hyperledger Fabric}\par
  \begin{itemize}
    \item Спроектировал и реализовал на Go вычислительный PIR-протокол на основе BGV
    с использованием Lattigo, интегрировав приватное чтение из мирового состояния
    непосредственно в чейнкод Hyperledger Fabric.
    \item Сформулировал модель угроз с честным, но любопытным противником (HBC) и
    проанализировал конфиденциальность запросов на основе IND-CPA-безопасности BGV,
    включая безопасность повторных запросов и наблюдаемую утечку информации.
    \item Построил экспериментальный стенд и оценил корректность, задержку, объём
    коммуникации, требования к хранению и накладные расходы развёртывания для
    нескольких наборов криптографических параметров.
  \end{itemize}

  \vspace{0.42em}

  \textbf{Агрегация с сохранением конфиденциальности для федеративного обучения с подкреплением}\par
  \begin{itemize}
    \item Участвовал в разработке TWAPPA-FRL, объединяющей доверительно-взвешенную
    агрегацию (TWA) и CKKS-агрегацию с сохранением конфиденциальности (PPA)
    для совместного обнаружения вторжений, в качестве автора с равным вкладом.
    \item Объединил взвешивание клиентов на основе доверия с зашифрованной
    CKKS-агрегацией обновлений федеративной модели и расшифрованием только
    агрегированного результата отдельным центром управления ключами,
    не участвующим в обучении.
    \item Разработал модель угроз и анализ безопасности, включая конфиденциальность
    индивидуальных обновлений, явно определённую утечку, а также предположения
    и ограничения путей выполнения с PPA.
  \end{itemize}

  \vspace{0.42em}

  \textbf{Верифицируемое многоключевое шифрование с возможностью поиска}\par
  \begin{itemize}
    \item Участвовал в разработке EVMK-SSE --- эффективной схемы многоключевого
    шифрования с возможностью поиска на основе OT и ослабленной OPRF,
    не зависящей от клиента.
    \item Проверил анализ безопасности, охватывающий конфиденциальность запросов,
    явную утечку шаблонов поиска и доступа, а также верифицируемость результатов
    при наличии вредоносного облачного сервера.
  \end{itemize}
}

\vspace{0.8em}

\cvrow{Опыт работы}{%
  \entry
    {\textbf{Стажёр по информационной безопасности} @
    \href{https://global.ptsecurity.com/en/}{Positive Technologies}}
    {авг. 2024--февр. 2025}
  \vspace{0.15em}

  \begin{itemize}
    \item Разработал протестированные инструменты на Python для команды MaxPatrol EDR,
    предназначенные для проверки путей Windows, сетевых конфигураций и политик
    контроля конечных устройств.
    \item Построил конвейер обработки вредоносного ПО на базе Airflow:
    сбор образцов из VX-Underground, хранение артефактов в S3/MinIO
    и сканирование с использованием YARA.
    \item Исследовал эксплуатацию и постэксплуатационную активность, связанную с
    CVE-2020-5902 и CVE-2023-38831, включая анализ EVTX, деобфускацию PowerShell,
    извлечение IoC и сопоставление техник с MITRE ATT\&CK.
    \item Развил полученный опыт в области анализа киберугроз в самостоятельный
    исследовательский проект MFCTIIF по интеграции нескольких источников CTI,
    выступив первым автором статьи.
  \end{itemize}

  \vspace{0.5em}

  \entry
    {\textbf{Инженер-программист} @ \href{https://academai.ru/}{AcademAI}}
    {авг. 2023--авг. 2024}
  \vspace{0.15em}

  \begin{itemize}
    \item Совместно разработал полнофункциональную AI-платформу для генерации
    структурированных и персонализированных образовательных курсов по темам
    и учебным целям, заданным пользователями.
    \item Реализовал REST API на Python и конвейер генерации на базе LangChain
    с использованием структурированных промптов и асинхронной оркестрации LLM.
    \item Разработал фронтенд и прикладной слой на Next.js, TypeScript, Prisma и Auth.js,
    включая аутентификацию и постоянное хранение данных курсов.
    \item Контейнеризировал и развернул платформу с использованием Docker и GitHub Actions.
  \end{itemize}
}

\vspace{0.8em}

\cvrow{Навыки}{%
  \renewcommand{\arraystretch}{1.22}
  \begin{tabularx}{\linewidth}
    {@{}>{\raggedright\arraybackslash}X
    @{\hspace{1.2em}}
    >{\raggedright\arraybackslash}X@{}}

    $\diamond$ \textbf{Go, Python, TypeScript:} уверенный уровень
    &
    $\diamond$ \textbf{C++, Rust:} базовый уровень
    \\

    $\diamond$ \textbf{Гомоморфное шифрование:} BGV, BFV, CKKS
    &
    $\diamond$ \textbf{Приватный поиск:} PIR/CPIR, шифрование с возможностью поиска
    \\

    $\diamond$ \textbf{Библиотеки ГШ:} Lattigo, Microsoft SEAL, OpenFHE, TenSEAL
    &
    $\diamond$ \textbf{Системы:} Hyperledger Fabric, Linux, Docker, Git
    \\

    $\diamond$ \textbf{Инженерия безопасности:} YARA, Chainsaw,
    анализ вредоносного ПО и журналов
    &
    $\diamond$ \textbf{Языки:} русский (родной),
    \href{https://certs.duolingo.com/751e0ef152565727b6d5fd909699ceaa}
        {английский (C1)},
    \href{https://admin.chinesetest.cn/queryScore.do?sid=c2tActR9Wb\%2704b861aefbdc5e9fd0b44afa717d08cbceb0720e582ce25f\%2724d7sWsQ5nVUQSwD}
        {китайский (HSK 3)}
  \end{tabularx}
}

\end{document}
